 \chapter{Conclusions}
The proposed Flow Planner is a generative hierarchical framework that addresses long-horizon decision-making problems by decoupling strategic reasoning from reactive control. This architecture divides planning into two distinct levels: a high-level global planner operating in the observation space to provide long-term strategic direction, and a low-level reactive action planner generating robust, short-horizon policies. This decomposition significantly reduces the complexity of the state space each flow model must learn. Consequently, the action planner can be trained efficiently on short-horizon tasks (e.g., 25 steps) while increasing model capacity (e.g., 64 hidden dimensions → 128) to avoid the sample inefficiency and challenges typically associated with end-to-end, long-horizon trajectory planning.

This hierarchical approach yields substantial computational and design benefits. Within an MPC loop, the computationally intensive global planner is invoked infrequently (e.g., once every 20 steps), leading to a low amortized inference time that makes the system practical for real-time applications. Furthermore, the framework's inherent modularity establishes a clear separation of concerns. This allows for independent optimization and debugging of each component: the global planner, the reactive planner, and the controller. This modularity facilitates the integration of hybrid systems that combine the strengths of both learned and classical planning methods.

While the empirical results in this dissertation demonstrated the superiority of Flow Matching for the Lunar Lander task, the framework's design does not advocate for one paradigm over the other. Instead, perhaps the most significant contribution is how the hierarchical structure provides a principled way to select the optimal generative model for different types of planning tasks. For problems that demand strategic exploration across the solution space, where fine-grained controllability and the generation of diverse, stochastic plans are necessary, diffusion models can be a superior paradigm. On the other hand, for tasks where raw sampling speed, deterministic precision, and efficient exploitation are paramount, Flow Matching is a lot more compelling. In this sense, the architecture brings the discussion full circle, grounding the choice between modern generative models in the classical reinforcement learning tension between exploration and exploitation.

This trade-off is not incidental but is a direct consequence of the models' fundamental formulations. Diffusion models learn a score function to reverse a fixed noising process, making their iterative denoising inference a natural fit for guided, stochastic exploration of the solution space. In contrast, Flow Matching learns a deterministic vector field to directly transport a prior to the data distribution. Its inference is a single ODE integration, which is inherently faster and produces smoother, more direct trajectories, making it the ideal choice for efficient, exploitative control. Ultimately, the choice between these powerful generative tools is a deliberate design decision, and the Flow Planner framework provides a principled architecture for leveraging the distinct strengths of both.


\section{Evaluation}
This work successfully achieved its primary objectives, culminating in the design and validation of a new class of generative planners based on the Flow Matching paradigm. The core technical contribution, presented in Chapter 3, was the novel reformulation of the foundational CFM framework to generate trajectories for robotic control, establishing a new class of deterministic, high-speed generative planners. This formulation was then incorporated into the Flow Planner architecture, a hierarchical framework that divides the planning problem into strategic and reactive components. This hierarchical structure proved essential for overcoming the scaling and long-horizon challenges inherent in monolithic, end-to-end planners.
The extensive empirical evaluation in Chapter 5 validated the framework's ability to generate high-quality trajectories that consistently matched the performance of an expert PPO agent. The direct comparison with the diffusion paradigm quantified the significant advantages of the flow-based approach in terms of inference speed and training stability, providing practical guidance for future research. Finally, this work released three comprehensive benchmark datasets along with their corresponding trained expert agents, enabling full reproducibility and providing valuable resources for future research.
\section{Future Work}
\subsection{Scaling Up to Different Environments}
The present work, developed within the Box2D environment, sets a solid foundation for Flow Planner, but its utility will be best demonstrated in more complex domains. Extending experimentation to Gymnasium environments such as Atari and MuJoCo \cite{todorov2012mujoco},  and potentially benchmarking on the latest OG-Bench tasks \cite{parkOGBenchBenchmarkingOffline2025}, would stress-test the planner’s adaptability. Pushing further into realistic robotic arms, using sim-to-real transfer techniques, would offer compelling evidence of robustness and effectiveness. Such scaling efforts would validate the approach’s applicability beyond controlled simulations and pave the way toward real-world deployment.

\subsection{Flow Optimization and other Applications}
There are many recent state-of-the-art advancements in optimizing flow for faster inference. For example, rectified flow is a technique that iteratively straightens the generated flow, enabling high-quality generation in a fraction of the ODE steps \cite{liuFlowStraightFast2022}. Similarly, Optimal Transport Flow (OT-CFM) could improve both training stability and inference speed by learning a more efficient transport map between the noise and data distributions \cite{tongImprovingGeneralizingFlowbased2024}. For tasks involving complex kinematics, such as 6-DoF manipulation, exploring Riemannian Flow Matching would allow the model to respect the underlying manifold geometry of the robot's configuration space (e.g., SO(3)), leading to more natural and physically plausible plans \cite{lipmanFlowMatchingGuide2024a}.

\subsection{Strengthening and Diversifying Conditioning Signals}
So far, conditioning in Flow Planner has been limited to start and end states, which constrains its expressive potential. A logical next step would be to implement reward-based conditioning, allowing the model to prioritize trajectories with higher expected returns, resulting in more consistent and optimal behavior. Additionally, integrating frameworks like ControlNet \cite{zhangAddingConditionalControl2023} could substantially enhance the conditioning mechanism by enabling external guidance signals (e.g. natural language task descriptions, environmental constraints, etc.) to influence trajectory generation. This enhancement would dramatically increase the controller's versatility across diverse tasks with varying requirements.

\subsection{Integration with Conventional Reinforcement Learning}
While Flow Planner has so far focused on imitation-style learning, its architecture is well-suited for integration with conventional reinforcement learning. As mentioned in Section \ref{sec:other_works} relevant works, Future work should explore incorporating methods such as Q-Score Matching (QSM) or Flow Q-Learning (FQL) to enable direct policy refinement \cite{psenkaLearningDiffusionModel2025b, parkFlowQLearning2025}. This would allow the planner to optimize behavior directly according to task-specific reward functions rather than purely reproducing trajectories.

\subsection{Towards a Hierarchical Reasoning Model}
The current Flow Planner architecture mirrors the dual-process "System 1" (fast, reactive) and "System 2" (slow, deliberative) thinking popularized by Kahneman. However, it lacks a crucial feedback loop from the low-level planner back to the high-level planner. A powerful future direction would be to implement an "iterative refinement loop," inspired by the recent Hierarchical Reasoning Model \cite{wangHierarchicalReasoningModel2025a} that demonstrated exceptional performance on the ARC-AGI benchmark\cite{cholletARCPrize20242025}. This approach would create a feedback loop between global and local planners, allowing the low-level planner to evaluate feasibility and projected costs before executing on the strategic “imaginations”. This is a mechanism shown to be critical for performance on complex, multi-step problems.